% Drop-in manuscript update generated from the formal questionnaire analysis.
% Formal dataset: U01--U10 only (N=10).
% U0 and U11 were incomplete. U12--U16 were synthetic supervisor-preview records
% and MUST NOT be included in formal human-subject results.

% ---------- Replace Section 5.1 ----------
\subsection{Participants and Overall Procedure}
\label{sec:eval_participants}

We analyzed ten usable participant records from the formal study.
Participants had character-animation or closely related performance-authoring
backgrounds spanning game/film animation, animation/VFX study, and professional
character animation. Seven participants reported a numeric duration of relevant
experience, ranging from 2 to 8 years ($M=4.57$, $SD=1.99$). Of the nine
participants who completed the software-background items, all reported prior
experience with Maya or another keyframe-animation package. We use animation
background and production experience only as descriptive participant
characteristics.

Participants first received an explanation of the semantic fields exposed by
the Semantic Performance Tag representation and completed a tutorial example.
The same participants then completed the end-to-end authoring tasks and the
text-only context-ablation ratings. Formal stimuli did not overlap across the
two studies.

% NOTE TO AUTHORS: fill session duration, compensation, and ethics reference from
% the approved study record; they are not recoverable from the questionnaire data.

% ---------- Replace the final paragraph of Section 5.2.1 ----------
Each participant completed two formal scenes with Direct Generation followed by
two scenes with Editable Semantic Performance Tags. The four authoring scenes
were distributed across participants, but the workflow order was fixed in the
collected sessions. Consequently, Study~A is analyzed as a paired first-use
comparison and any workflow difference must be interpreted with a possible
practice/order effect rather than as an order-independent causal estimate.
Direct low-level Maya curve editing was excluded from the formal task.

% ---------- Replace Section 5.2.3 ----------
\subsubsection{Subjective Measures}
After each formal authoring task, participants rated (1) how well the final
performance matched their intended reading of the scene and (2) satisfaction
with the final performance on seven-point scales. After each two-task workflow
block, participants completed the six raw NASA-TLX dimensions. They also rated
nine workflow statements individually on seven-point agreement scales,
covering expression of acting direction, local revision while preserving
accepted choices, revision clarity, control of the final acting direction,
speaker--listener coordination, movement between visual judgment and semantic
revision, creative ownership, exploration of alternatives, and whether the
result justified the invested effort. We report these items individually rather
than combine them into an unvalidated composite score.

After the Editable Semantic Performance Tag block, participants additionally
rated three diagnostic items concerning whether the tags matched the level at
which they wanted to edit performance, whether phrase-level acting
interpretations were useful, and whether the dual-character layout supported
speaker--listener coordination. These diagnostic items were not used in the
between-workflow statistical comparison.

% ---------- Add at the end of Section 5.2 ----------
\subsubsection{Study A Analysis}
For the two task-level outcomes, we first averaged each participant's two tasks
within each workflow. The nine workflow items and six NASA-TLX dimensions were
analyzed as paired participant-level observations. Because the sample is small
and the primary questionnaire outcomes are ordinal, we use two-sided Wilcoxon
signed-rank tests. We report rank-biserial correlations as paired effect sizes
and apply Holm correction within the task-outcome, workflow-item, and NASA-TLX
families. Positive rank-biserial values indicate higher ratings in the Editable
Semantic Performance Tag condition; for NASA-TLX, higher values indicate
greater workload except for the performance dimension, where 0 denotes perfect
performance and 100 denotes failure.

% ---------- Replace Section 5.3.3 ----------
\subsubsection{Study B Analysis}
The primary outcome is overall contextual appropriateness; acting intention,
affect, attention/gaze, and temporal placement are secondary outcomes. For each
criterion, we averaged the six Dialogue-Only ratings and the six Full-Context
ratings within participant and compared the paired participant-level values
with two-sided Wilcoxon signed-rank tests. We report rank-biserial correlations
and Holm-adjusted $p$ values across the five criteria. Starting-plan choices are
reported descriptively after decoding the counterbalanced anonymous Plan A/Plan
B assignments back to Dialogue-Only and Full-Context; ``No preference'',
missing, and ambiguous responses are retained separately.

% ---------- Add Section 5.4 ----------
\subsection{Results}
\label{sec:results}

\subsubsection{Study A: Authoring Workflow}
Editable Semantic Performance Tags produced numerically higher final-performance
ratings, but the differences were not statistically reliable in this sample.
Mean match to the participant's intended reading increased from
$5.10$ ($SD=0.97$) with Direct Generation to $5.65$ ($SD=0.82$) with editable
tags ($W=9.0$, $p=.117$, Holm-adjusted $p=.219$, $r_{rb}=.60$).
Final-performance satisfaction similarly increased from $5.15$ ($SD=1.06$) to
$5.75$ ($SD=0.89$; $W=4.0$, $p=.109$, Holm-adjusted $p=.219$,
$r_{rb}=.71$). Figure~\ref{fig:study_a_task} summarizes these task-level
outcomes.

\begin{figure}[t]
  \centering
  \includegraphics[width=\columnwidth]{fig_studya_task_outcomes.png}
  \caption{Study A final-performance ratings. Bars show participant-level means
  with standard-deviation error bars. The formal analysis includes ten usable
  participant records.}
  \label{fig:study_a_task}
\end{figure}

Across the nine workflow items, the largest directional difference concerned
revision clarity: Direct Generation $M=4.89$ ($SD=1.54$) versus editable tags
$M=6.22$ ($SD=0.67$), $W=0$, $p=.063$, Holm-adjusted $p=.563$,
$r_{rb}=1.00$ ($n=9$ paired responses). Dyadic coordination also increased
directionally from $M=5.00$ ($SD=1.49$) to $M=5.80$ ($SD=0.92$),
but this difference was not significant ($p=.250$; Holm-adjusted $p=1.00$).
No workflow item remained significant after family-wise correction
(Table~\ref{tab:workflow_results}).

\begin{table*}[t]
\centering
\caption{Study A workflow-item ratings. Positive $r_{rb}$ favors Editable
Semantic Performance Tags. Holm correction is applied across the nine items.}
\label{tab:workflow_results}
\small
\begin{tabular}{lrrrrrr}
\toprule
Item & $n$ & Direct $M$ & Tags $M$ & $W$ & $p_{\mathrm{Holm}}$ & $r_{rb}$\\
\midrule
Express intended acting direction & 10 & 5.50 & 5.80 & 2.0 & 1.000 & .60\\
Local revision while preserving choices & 10 & 5.00 & 5.50 & 13.0 & 1.000 & .28\\
Understand what was being changed & 9 & 4.89 & 6.22 & 0.0 & .563 & 1.00\\
Control of final acting direction & 10 & 5.20 & 5.60 & 4.0 & 1.000 & .47\\
Coordinate speaking/listening behavior & 10 & 5.00 & 5.80 & 2.0 & 1.000 & .73\\
Move between visual judgment and revision & 9 & 5.22 & 5.00 & 2.5 & 1.000 & -.50\\
Creative ownership & 9 & 4.78 & 5.33 & 2.5 & 1.000 & .67\\
Explore alternative performances & 10 & 6.30 & 6.40 & 2.0 & 1.000 & .33\\
Quality worth time/effort & 8 & 6.13 & 6.13 & 5.0 & 1.000 & .00\\
\bottomrule
\end{tabular}
\end{table*}

Raw NASA-TLX scores did not show a statistically reliable workload difference
between workflows on any dimension after Holm correction. Mental demand was
$39.3$ ($SD=18.1$) for Direct Generation and $41.5$ ($SD=19.2$) for editable
tags; effort was $37.0$ ($SD=21.6$) and $40.3$ ($SD=24.2$), respectively.
The NASA-TLX performance dimension (0=perfect, 100=failure) was numerically
higher for editable tags ($M=55.2$) than Direct Generation ($M=43.0$), but the
unadjusted comparison did not cross the conventional threshold
($p=.055$; Holm-adjusted $p=.328$). We therefore find no evidence that the
editable representation reduced workload in this first-use study.

The three tag-specific diagnostics were positive but should be read as
descriptive. Participants rated the match between Semantic Performance Tags and
their desired level of editing at $M=5.00$ ($SD=1.05$), the usefulness of
phrase-level acting interpretation at $M=5.30$ ($SD=1.34$), and the
dual-character layout for coordinating speaker and listener performance at
$M=6.20$ ($SD=0.63$).

\subsubsection{Study B: Context Ablation}
Full-Context proposals received higher overall scene-fit ratings than
Dialogue-Only proposals. Participant-level means increased from
$5.22$ ($SD=0.90$) to $5.58$ ($SD=1.02$), with all non-zero participant
differences favoring Full-Context ($W=0$, $p=.0078$, Holm-adjusted $p=.039$,
$r_{rb}=1.00$). This was the only criterion that remained significant after
Holm correction. Intention ($5.20$ vs.\ $5.43$), affect ($5.05$ vs.\ $5.08$),
attention/gaze ($5.42$ vs.\ $5.48$), and temporal placement ($5.45$ vs.\
$5.59$) did not show corrected significant differences
(Table~\ref{tab:context_results}). Thus, in the tested setup, added context
changed the holistic appropriateness of a proposed performance more clearly
than any single exposed semantic channel.

\begin{table}[t]
\centering
\caption{Study B participant-level ratings for Dialogue-Only (D) and
Full-Context (F) proposals. Holm correction is across the five criteria.}
\label{tab:context_results}
\small
\begin{tabular}{lrrrr}
\toprule
Criterion & D $M$ & F $M$ & $p_{\mathrm{Holm}}$ & $r_{rb}$\\
\midrule
Overall scene fit & 5.22 & 5.58 & .039 & 1.00\\
Character intentions & 5.20 & 5.43 & 1.000 & .38\\
Facial affect & 5.05 & 5.08 & 1.000 & .19\\
Attention/gaze & 5.42 & 5.48 & 1.000 & .21\\
Change timing & 5.45 & 5.59 & .375 & .86\\
\bottomrule
\end{tabular}
\end{table}

Across 60 scene-level preference questions, Full-Context was selected 27 times,
Dialogue-Only 15 times, and ``No preference'' 11 times; six responses were
missing and one response was ambiguous. Among the 42 determinate choices for
one of the two proposals, Full-Context therefore accounted for 64.3\%.
We treat this as descriptive support for the appropriateness result rather than
as an independent inferential test.

% ---------- Replace Discussion 6.1 ----------
\subsection{When Does Context Change a Performance Interpretation?}
Study~B provides evidence that narrative and interpersonal context matters at
the level of the performance as a whole. Full-Context proposals improved
overall scene-fit ratings after correction, while the individual ratings for
intention, affect, gaze, and temporal placement did not independently reach
corrected significance. One interpretation is that context changes the
coherence among several plausible choices rather than simply moving a single
facial channel in a predictable direction. This is consistent with the design
premise that contextual information constrains an acting space without
establishing one uniquely correct performance. The starting-plan preferences
point in the same direction, but the presence of no-preference and missing
responses cautions against treating Full-Context as uniformly preferable.

% ---------- Replace Discussion 6.2 ----------
\subsection{Performance Interpretation as an Authoring Object}
Study~A does not support a strong claim that exposing Semantic Performance Tags
improves control, workload, or final quality. Final-performance ratings moved
in a favorable direction and revision clarity showed the largest paired effect,
but none of the workflow comparisons survived multiple-comparison correction
in this ten-participant sample. The tag-specific diagnostics nevertheless
suggest that participants generally understood the representation and
especially valued the dual-character coordination view. These findings are
better read as evidence that the semantic layer was usable as an authoring
surface than as evidence of superiority over prompt-level regeneration.

The absence of a workload reduction is also informative. Editable semantic
control introduces an additional object that must be inspected and revised.
The present data therefore fit a control--effort trade-off more closely than an
efficiency narrative. A larger counterbalanced study would be required to
separate the value of semantic locality from practice and order effects.

% ---------- Replace Discussion 6.4 ----------
\subsection{Interpretation Visibility Without Treating Reasons as Ground Truth}
Phrase-level acting interpretation received a mean diagnostic rating of 5.30
out of 7, indicating moderate-to-positive perceived usefulness, but lower than
the 6.20 rating for the dual-character layout. This supports keeping the
interpretation available as lightweight provenance rather than making
explanation the center of the workflow. The representation communicates the
system's proposed acting reading so that it can be accepted or contested; it
does not establish that the proposal is correct or expose hidden model
reasoning.

% ---------- Add to Limitations 6.5 ----------
The final authoring sample was small ($N=10$), and the collected Study~A
sessions used a fixed Direct-Generation-then-editable-tags order. Practice,
familiarity with the scenes and interface, or fatigue may therefore contribute
to paired workflow differences. The present quantitative analysis is also
stronger for RQ1 and RQ2 than for RQ3: the questionnaire data characterize
ratings and preferences, but claims about which generated semantic decisions
were accepted, modified, or removed require the original-to-final plan diffs
and session logs to be analyzed separately. We therefore do not infer
acceptance or rejection rates from the questionnaire alone.

% ---------- Replace the result placeholder in Conclusion ----------
In the context ablation, Full-Context proposals were rated higher on overall
scene fit than Dialogue-Only proposals, while individual affect, gaze, intention,
and timing ratings did not show corrected differences. In the first-use
authoring comparison, editable Semantic Performance Tags produced directional
improvements in final-performance ratings and revision clarity but no corrected
significant workflow or workload advantage, supporting a cautious view of the
representation as an inspectable authoring surface rather than a demonstrated
replacement for prompt-level regeneration.
